%! Rates of Change in Polar Functions — Unit 7, Lesson 7.8 — Student Packet Cover Sheet
\documentclass[10pt]{article}
\usepackage{precalculus-article}
\usepackage{precalculus-boxes}
\usepackage{ltablex}
\keepXColumns

\begin{document}

% Full-bleed navy banner
\begin{tikzpicture}[remember picture, overlay]
  \fill[navy] (current page.north west)
    rectangle ([yshift=-0.9in]current page.north east);
\end{tikzpicture}
\vspace{-0.8in}

\begin{center}
  {\color{white}\textbf{\LARGE Precalculus}}\\[4pt]
  {\color{skymid}\large\textbf{Unit 7: Analytic Trigonometry and Polar Coordinates}}\\[6pt]
  {\color{white}\textbf{\large Lesson 7.8 \quad Rates of Change in Polar Functions}}
\end{center}

\vspace{0.22in}
\namedateperiod
\vspace{0.18in}

\begin{learningtargetbox}
\small
\begin{enumerate}[leftmargin=1.5em, itemsep=5pt]
  \item I can compute the average rate of change of a polar function $r = f(\theta)$
    over an interval $[\theta_1, \theta_2]$ using the formula
    $\dfrac{f(\theta_2)-f(\theta_1)}{\theta_2-\theta_1}$, and interpret the result
    as how quickly the radius changes per radian.
    \hfill\textit{(LO 3.15.A, EK 3.15.A.1)}
  \item I can classify an interval of a polar function as positive and increasing,
    positive and decreasing, negative and increasing, or negative and decreasing,
    and describe whether the curve is moving toward or away from the origin.
    \hfill\textit{(LO 3.15.A, EK 3.15.A.2)}
  \item I can explain how the sign (positive/negative) of $r$ combined with whether
    $r$ is increasing or decreasing determines the motion of the polar curve in the
    plane, including the opposite-ray behavior when $r < 0$.
    \hfill\textit{(LO 3.15.A, EK 3.15.A.3)}
\end{enumerate}
\end{learningtargetbox}

\vspace{0.14in}

\begin{tocbox}
\small
\renewcommand{\arraystretch}{1.5}
\begin{tabularx}{\linewidth}{@{} c l X r @{}}
\rowcolor{sky}
\textbf{\#} & \textbf{Component} & \textbf{Description} & \textbf{Score} \\
\midrule
1 & Warm-Up        & Spiral review: average rate of change, evaluating polar functions, identifying increasing/decreasing intervals & \blank{1.2cm} \\
2 & Guided Notes   & ARC formula and interpretation; four-case sign/direction classification; reading intervals from a table & \blank{1.2cm} \\
3 & Group Activity & Radar context: classify intervals, compute ARCs, interpret rate of change (tiered R/A/E) & \blank{1.2cm} \\
4 & Exit Ticket    & Classify two intervals; compute ARC; interpret negative rate of change & \blank{1.2cm} \\
5 & Homework       & Classify intervals for $r=3\sin\theta$; compute ARCs; interpret; AP-style MC & \blank{1.2cm} \\
\midrule
  & \multicolumn{2}{r}{\textbf{Total}} & \blank{1.2cm} \\
\end{tabularx}
\end{tocbox}

\end{document}
